\chapter{Revisão Bibliográfica}\label{cap:fund_teorica}

\section{Postura sentada e avaliação postural}
\label{sec:rev-ergonomia}

Os DORT afetam músculos, tendões, articulações e nervos. Entre seus fatores de
risco estão a postura inadequada, a permanência prolongada em posição estática
e a ausência de pausas, agravados por fatores psicossociais
\cite{mohammadian2025}. O risco postural costuma ser avaliado por métodos
observacionais: o RULA \cite{mcatamney1993} pontua faixas de \textbf{ângulo de
inclinação do tronco e do pescoço} --- o pescoço é penalizado em faixas de
10\degree{} de flexão ---, e a norma ISO 11226 \cite{iso11226} combina o ângulo
com o \textbf{tempo} de manutenção da postura. Ângulo do tronco e persistência
no tempo fundamentam, neste trabalho, a inclinação medida por IMU e a espera
antes do alerta; os limiares do sistema, porém, são desvios em relação à
postura de referência de cada usuário, e não ângulos absolutos de norma.

\section{Sensores para monitoramento postural}
\label{sec:rev-sensores}

\textbf{Resistores sensíveis à força (FSR).} O FSR é um filme polimérico cuja
resistência cai à medida que a força aplicada aumenta. O FSR402 tem área ativa
de 12{,}7\,mm de diâmetro, faixa de cerca de 0{,}2 a 20\,N, curva
força--resistência fortemente não linear, repetibilidade de $\pm$6\,\% entre
peças e deriva sob carga constante (\emph{creep}) inferior a 5\,\% por década
de tempo \cite{interlink_datasheet}. O \emph{creep}, de origem viscoelástica, é
a limitação mais relevante para quem permanece sentado por longos períodos
\cite{paredesmadrid2017}, e o fabricante adverte que o FSR não é uma célula de
carga: a exatidão em força varia de $\pm$6\,\% a $\pm$50\,\% conforme a
montagem \cite{interlink_guia}. O FSR serve, portanto, como medida
\textbf{relativa} de distribuição de carga, comparando sensores entre si ou com
uma linha de base própria.

\textbf{Acelerômetros e IMUs.} Em repouso, o acelerômetro mede o vetor
gravidade no referencial do sensor, do qual se obtêm os ângulos de arfagem
(\emph{pitch}) e rolagem (\emph{roll}). A estimativa supõe aceleração linear
desprezível, não observa a guinada (\emph{yaw}) e é instável perto de
$\pm$90\degree{} de arfagem \cite{pedley2013}; movimentos dinâmicos exigem
combiná-lo ao giroscópio. Como a postura sentada é quase estática, este
trabalho usa apenas o acelerômetro. O ângulo absoluto, porém, depende de como
o sensor assenta nas costas a cada uso, o que motiva a calibração por usuário.

\section{Trabalhos relacionados}
\label{sec:rev-relacionados}

As primeiras cadeiras instrumentadas usavam mapas de pressão de alta
resolução \cite{tan2001}; as seguintes passaram a poucos sensores discretos.
\citeonline{zemp2016} combinaram 16 FSRs a um módulo inercial no encosto, que
mede a inclinação da cadeira, e não a do usuário, e classificaram sete
posições. Nesses trabalhos a saída é a classificação da postura, sem retorno ao
usuário. O LifeChair, almofada com sensores de pressão, vibração e aplicativo,
acrescenta o retorno e observa que o índice de massa corporal altera o padrão
de pressão \cite{bourahmoune2019}. Do lado vestível, um sensor no tronco que
vibra quando este se curva reduziu a flexão do pescoço e a cifose torácica de
21 adultos durante uma hora de digitação \cite{kuo2021}. A vibração não ocupa
a visão, já dedicada à tela, mas é percebida como o retorno mais intrusivo
\cite{haller2011}, e a evidência de efeito sustentado é fraca: em ensaio com
trabalhadores, o retorno sobre a postura lombar não diferiu do grupo controle
\cite{ribeiro2014}. O alerta deve, assim, disparar só depois que o desvio
persiste, e o efeito em longo prazo não pode ser presumido a partir de testes
curtos.

O BackWatch \cite{ferreira2025}, pôster da FEBRACE 2025, mede com dois
sensores ultrassônicos no encosto o afastamento do usuário e alerta por display
LCD acima de um limiar fixo; é montado na cadeira, e não vestível. Na revisão
sistemática de \citeonline{kappattanavar2021}, 13 dos 14 estudos usavam
sensores de pressão e nenhum os combinava a um IMU no corpo; os autores
recomendaram incorporar IMUs para capturar o movimento da coluna.

O sistema proposto ocupa essa lacuna em quatro aspectos: \textbf{funde} a
distribuição de peso na cadeira com a orientação do tronco superior, que não
toca a cadeira, de modo que cada canal cobre casos que o outro confunde
(reclinar inclina o tronco sem ser prejudicial; escorregar descarrega o apoio
traseiro do assento com pouca inclinação); fornece \textbf{retorno háptico} em
tempo real no corpo; usa \textbf{calibração por usuário}, em vez de limiares
absolutos ou modelos treinados em outras pessoas; e faz todo o
\textbf{processamento embarcado}, sem aplicativo nem nuvem.
